% Chapter 3, section 1: Work queue
% Write in plain LaTeX. Code of a guide step: \guide{<step>}; commands: \begin{shell}...\end{shell};
% a screenshot: \screenshot{file}{caption}; a gap in the draft: \placeholder{...}.
% See book/latex/queen.sty for the full list.
%
% Draft by Claude, 2026-10-24. The code is cut from examples/tutorials/js/book-use-cases.mjs, which
% checks every result quoted here against a 2.7.0 node.

\section{Work queue}

Many workers share one stream of jobs, and each job is done by exactly one of them: resizing
photos, sending emails, rendering invoices. This is the work queue, also called competing
consumers.

In Queen it is a consumer with no consumer group. Each message goes to one worker. Work runs in
parallel across partitions: a partition is held by one worker at a time, so here every user's
photos are a partition, and three workers share four users.

\queenfig{work-queue}{Competing consumers: each job once, one user's partition per worker at a time.}

\guide{uc-work}

All twelve photos were resized once each, shared out as 6, 3 and 3 between the three workers.

\why{order still holds inside each user's photos, and a slow user holds up only its own
partition. To go faster, add workers, up to the number of partitions that have work.}

One thing to know: with a single partition the workers would take turns, one at a time. Give the
jobs a partition per user, per customer or per tenant.
